\documentclass[a4paper]{article}
\usepackage{amsmath}
\usepackage[utf8]{inputenc}
\usepackage[english]{babel}

\usepackage[style=ieee]{biblatex}
\addbibresource{channel.bib}


\title{Channel models for terrestrial radio and acoustic underwater links}
\author{Bálint Áron Üveges}
\date{\today}

\begin{document}
\maketitle

\section{Castalia's path loss model}
Castalia utilizes the Log-distance path loss model:
\begin{equation}
    PL_{d} = PL_{d0} + 10.0 \cdot \gamma \cdot log_{10} \left( \frac{d}{d_0} \right) +  \mathcal{N}(0,\sigma) + \frac{\mathcal{N}(0,\sigma_{b})}{2}
\label{eq:log_d_pl}
\end{equation}
In \eqref{eq:log_d_pl}, $PL_{d0}$ is the path loss at the reference distance $d_{0}$; $d$ is the length of the path; $\gamma$ is the path loss exponent; $\mathcal{N}$ is the Gaussian-distribution, $\sigma$ is the standard variance of shadow fading; whereas $\sigma_{b}$ is the standard variance of bidirectional path loss jitter.

\section{Underwater acoustic channel path loss model}
According to \cite{Stojanovic_2007}, the narrowband 

\section{Castalia's parameters, bits and pieces}
defaultChannel:
\begin{itemize}
    \item \texttt{signalDeliveryThreshold} -- threshold in dBm above which, wireless channel module is delivering signal messages to radio modules of individual nodes.
\end{itemize}

radio:
\begin{itemize}
    \item \texttt{carrierFreq} -- The carrier frequency in MHz.
    \item \texttt{rssiIntegrationTime} -- rssiIntegrationTime = symbolsForRSSI * RXmode->bitsPerSymbol / RXmode->datarate
    \item \texttt{timeToTxPacket} -- Provided by popAndSendToWirelessChannel(), calculated as: double txTime = ((double)(end->getByteLength() * 8.0f)) / RXmode->datarate; (in seconds)
    \item \texttt{parameter file} -- Name, dataRate(kbps), modulationType, bitsPerSymbol, bandwidth(MHz), noiseBandwidth(KHz), noiseFloor(dBm), sensitivity(dBm), powerConsumed(mW)
    \item \texttt{SNR2BER} -- method used to calculate BER based on SNR and the actual modulation.
\end{itemize}

\section{AquaSim's parameters, bits and pieces}



\printbibliography
\end{document}
